\documentclass[12pt,a4paper]{article}
\usepackage[utf8]{inputenc}
\usepackage[T1]{fontenc}
\usepackage[brazil]{babel}   % requer o pacote de idioma portugues do babel (presente no Overleaf e em instalacoes TeX Live completas)
\IfFileExists{lmodern.sty}{\usepackage{lmodern}}{} % fonte opcional; ignora se ausente
\usepackage[margin=2.5cm]{geometry}
\usepackage{booktabs}
\usepackage{tabularx}
\usepackage{array}
\usepackage{ragged2e}
\usepackage[table]{xcolor}
\usepackage{titlesec}
\usepackage{enumitem}
\usepackage{setspace}
\usepackage[hidelinks]{hyperref}

\definecolor{accent}{HTML}{2E5A2E}
\definecolor{grey}{HTML}{555555}
\newcommand{\rs}{R\$}
\setlength{\parindent}{0pt}
\setlength{\parskip}{6pt}
\onehalfspacing
\titleformat{\section}{\color{accent}\large\bfseries}{\thesection.}{0.5em}{}
\titleformat{\subsection}{\bfseries}{\thesubsection}{0.5em}{}
\newcolumntype{Y}{>{\centering\arraybackslash}X}
\renewcommand{\arraystretch}{1.2}
\setlist[itemize]{leftmargin=1.4em,itemsep=2pt,topsep=2pt}
\setlist[enumerate]{leftmargin=1.6em,itemsep=2pt,topsep=2pt}

\begin{document}

% ---------- CAPA ----------
\begin{titlepage}
\centering
\vspace*{1cm}
{\color{grey}\large [Instituição de Ensino]\par}
\vspace{0.3cm}
{\color{grey}\large Disciplina: Engenharia Econômica\par}
\vfill
{\bfseries\Huge Plataforma Inteligente Integrada de Monitoramento Rural\par}
\vspace{0.6cm}
{\color{accent}\Large\itshape IA $\cdot$ Drones $\cdot$ IoT $\cdot$ Visão Computacional\par}
\vspace{0.4cm}
{\large Relatório de Viabilidade Econômica\par}
\vfill
{\large Equipe: [preencher]\par}
\vspace{0.2cm}
{\color{grey}\large Produto (nome a definir)\par}
\vfill
{\large Goiânia, GO --- 2026\par}
\end{titlepage}

% ---------- CORPO ----------
{\bfseries\LARGE Plataforma Inteligente Integrada de Monitoramento Rural\par}
\vspace{2pt}
{\color{grey}\itshape IA $\cdot$ Drones $\cdot$ IoT $\cdot$ Visão Computacional --- Relatório de Viabilidade (Engenharia Econômica)\par}
\vspace{2pt}
{\color{grey}\small Equipe: [preencher] $\cdot$ Data: [preencher]/2026 $\cdot$ Produto (nome a definir)\par}

\vspace{6pt}
\noindent\colorbox{accent!8}{\parbox{\dimexpr\textwidth-2\fboxsep}{\small\itshape Os valores financeiros deste relatório são premissas de trabalho conferidas matematicamente, não dados de mercado verificados. A versão final deve substituí-los pelas cotações e preços de concorrentes coletados (Seção 4). A planilha \texttt{modelo\_financeiro.xlsx} recalcula tudo automaticamente ao editar as premissas.}}

\subsection*{Resumo}
Este relatório avalia a viabilidade econômica de uma plataforma B2B de monitoramento rural que integra drones, IoT e visão computacional com IA, voltada a médias e grandes fazendas de grãos. Comparam-se três modelos de receita (SaaS, DaaS e Híbrido) pelos indicadores de Engenharia Econômica --- VPL, TIR, payback simples e descontado, IBC e ROIA --- usando a Selic vigente como base da TMA. O modelo Híbrido (taxa de implantação + assinatura por hectare) é o recomendado, com VPL positivo e payback descontado de cerca de três anos. A análise de sensibilidade mostra que a viabilidade depende criticamente da velocidade de adoção de hectares.

\section{Contextualização}
O agronegócio brasileiro opera sob margens apertadas, exigência crescente de rastreabilidade e janelas de decisão curtas --- pragas, doenças e estresse hídrico comprometem a safra em poucos dias. Simultaneamente, três habilitadores amadureceram: conectividade rural em forte expansão (parcela relevante da área já com 4G/5G), queda no custo de drones e sensores, e IA/visão computacional com precisão operacional. O resultado é uma transição para ``fazendas inteligentes''. A demanda já está validada por concorrentes ativos (Perfect Flight, com dezenas de milhões de hectares digitalizados; SIMA, com milhões de hectares monitorados) --- o risco não é de mercado, é de execução comercial. O papel da análise econômica é duplo: provar que o investimento \textbf{se paga}, e em quanto tempo, tanto para o produtor quanto para o negócio.

\section{Problema (a dor do cliente)}
Médias e grandes fazendas de grãos perdem produtividade por \textbf{detecção tardia} de pragas, doenças, falhas de plantio e estresse hídrico, e gastam mais que o necessário com insumos por aplicação uniforme. Hoje resolvem com inspeção manual amostral e \textbf{ferramentas fragmentadas} (um app de satélite, outro fornecedor de drone, planilha de sensores, ERP à parte --- sem integração). As soluções atuais falham por: (i) fragmentação dos dados; (ii) pressuposto de conectividade estável, que quebra no campo; (iii) entrega de ``mapas'', não de \textbf{prescrições acionáveis}; e (iv) retorno não demonstrado em \rs, o que trava a adoção. A consequência prática é perda financeira recorrente que o produtor sequer mede com precisão.

\section{Proposta de solução}
Uma plataforma única que conecta \textbf{drones, sensores IoT e visão computacional com IA} e transforma imagens e leituras em \textbf{decisões acionáveis}: alertas priorizados de praga/doença, mapas de prescrição de taxa variável (VRT) exportáveis para o maquinário e indicadores de produtividade por talhão --- funcionando mesmo com conectividade ruim (operação \textit{offline-first}).

\textbf{Diferenciais defensáveis:} integração ponta-a-ponta (drone + IoT + satélite + ERP) em uma camada só; tradução em prescrição e em \rs, não apenas visualização; resiliência de conectividade; e base de dados proprietária que melhora os modelos a cada safra (efeito de rede de dados).

\section{Metodologia da solução}
\subsection{3.1 Técnico --- arquitetura}
A arquitetura é organizada em camadas, o que permite um MVP enxuto e escala sem reescrever o núcleo. A separação entre borda (campo) e nuvem é a decisão central: tarefas leves e o buffer offline ficam num \textit{gateway} na fazenda; o processamento pesado (ortomosaico, treino de modelos) fica na nuvem, com custo elástico.

\begin{table}[h]\small\centering
\begin{tabularx}{\textwidth}{>{\bfseries}l X X}
\toprule
\rowcolor{accent!12} Camada & Função & Racional \\
\midrule
Captação (Edge) & Drone (RGB/multiespectral) + IoT + satélite & Redundância de fonte de dado \\
Borda (Gateway) & Buffer offline, pré-processamento, inferência leve & Resolve conectividade; reduz custo de nuvem \\
Processamento (Nuvem) & Ortomosaico, treino, visão computacional pesada & Escala elástica de GPU \\
Dados & Geoespacial + séries temporais + imagens & Cada dado no banco adequado \\
Aplicação & Dashboard, alertas, mapas de prescrição & Onde o dado vira decisão \\
Integração & Exportar prescrição (ISOBUS/ISO-XML) e ERP & O ``último metro'' do valor \\
\bottomrule
\end{tabularx}
\end{table}

\textit{Faseamento:} MVP (uma fonte + um caso de IA + dashboard + 1 fazenda-piloto) $\to$ V1 (IoT, satélite e mapa de prescrição) $\to$ V2 (modelos preditivos, integrações com maquinário, base multi-fazenda).

\subsection{3.2 Financeiro --- viabilidade}
\textbf{Indicadores utilizados (e o que cada um mede):}

\begin{table}[h]\small\centering
\begin{tabularx}{\textwidth}{>{\bfseries}l X l}
\toprule
\rowcolor{accent!12} Indicador & O que mede & Regra de decisão \\
\midrule
TMA & Custo de oportunidade (melhor aplicação de baixo risco) & Taxa de desconto do VPL \\
VPL & Riqueza criada hoje, já descontada a TMA & Viável se VPL > 0 \\
TIR & Taxa que zera o VPL (rentabilidade intrínseca) & Atrativo se TIR > TMA \\
Payback simples & Tempo para recuperar o capital (sem desconto) & Quanto menor, melhor \\
Payback descontado & Tempo para recuperar pagando a TMA & Início real da criação de valor \\
IBC & Valor presente dos benefícios $\div$ investimento & Cria valor se IBC > 1 \\
ROIA & Ganho percentual real acima da TMA & Quanto maior, melhor \\
\bottomrule
\end{tabularx}
\end{table}

\textbf{TMA.} A referência atual é a \textbf{Selic = 14,50\% a.a. (Copom, jun/2026)}. Usamos três níveis: \textbf{piso 14,5\%} (custo de oportunidade puro), \textbf{recomendada 20\%} (Selic + prêmio de risco de venture) e \textbf{estresse 30\%}.

\textbf{Premissas centrais (ilustrativas --- recalibrar com dados reais):}

\begin{table}[h]\small\centering
\begin{tabularx}{\textwidth}{X r}
\toprule
\rowcolor{accent!12} Premissa & Valor \\
\midrule
Investimento inicial (Ano 0) & \rs~380.000 \\
Horizonte de análise & 5 anos \\
Curva de adoção (hectares ativos) & 25 mil $\to$ 550 mil ha \\
Fazenda média (médias/grandes grãos) & 3.000 ha/cliente \\
Preço SaaS (modelo recomendado) & \rs~22/ha/ano + implantação \\
\bottomrule
\end{tabularx}
\end{table}

\textbf{Comparação dos três modelos de receita} (à TMA recomendada de 20\% a.a.):

\begin{table}[h]\small\centering
\begin{tabularx}{\textwidth}{>{\bfseries}l Y Y Y}
\toprule
\rowcolor{accent!12} Indicador & A) SaaS & B) DaaS & C) Híbrido \\
\midrule
Investimento Ano 0 & \rs~380.000 & \rs~380.000 & \rs~380.000 \\
VPL (TMA 20\%) & \rs~2.788.869 & \rs~431.690 & \textbf{\rs~2.937.891} \\
VPL (TMA 14,5\%) & \rs~3.581.113 & \rs~843.042 & \rs~3.743.110 \\
VPL (TMA 30\%) & \rs~1.774.627 & $-$\rs~67.162 & \rs~1.903.646 \\
TIR & 79,9\% & 28,3\% & \textbf{85,6\%} \\
Payback simples & 2,91 anos & 4,21 anos & \textbf{2,70 anos} \\
Payback descontado & 3,20 anos & 4,70 anos & \textbf{3,07 anos} \\
IBC & 8,34 & 2,14 & 8,73 \\
ROIA & 52,8\% a.a. & 16,4\% a.a. & 54,2\% a.a. \\
\bottomrule
\end{tabularx}
\end{table}

\textbf{Leitura} (decisão entre alternativas excludentes --- método do Módulo 4 do material):
\begin{itemize}
\item \textbf{C (Híbrido)} é o melhor: a taxa de implantação adianta caixa e dá o menor payback com o maior VPL. Recomendado como modelo principal.
\item \textbf{A (SaaS asset-light)} é alternativa enxuta, simples de escalar, quase tão boa quanto C.
\item \textbf{B (DaaS frota própria)} parece atraente (ticket 3$\times$ maior), mas é intensivo em capital (46 unidades no Ano 5) e de margem fina (serviço de campo). Resultado: VPL muito menor, payback de $\approx$4,7 anos e valor negativo a 30\% de TMA. É a alternativa a \textbf{rejeitar} (ou oferecer só como serviço premium pago pelo cliente).
\end{itemize}

\textbf{Fluxo de caixa do modelo recomendado (C --- Híbrido):}

\begin{table}[h]\small\centering
\begin{tabularx}{\textwidth}{c Y Y Y Y}
\toprule
\rowcolor{accent!12} Ano & Receita (\rs) & Custo total (\rs) & Fluxo líquido (\rs) & Fluxo acum. (\rs) \\
\midrule
0 & --- & --- & $-$380.000 & $-$380.000 \\
1 & 716.667 & 1.115.000 & $-$398.333 & $-$778.333 \\
2 & 1.840.000 & 1.752.000 & 88.000 & $-$690.333 \\
3 & 4.120.000 & 3.136.000 & 984.000 & 293.667 \\
4 & 8.106.667 & 5.632.000 & 2.474.667 & 2.768.333 \\
5 & 13.633.333 & 9.090.000 & 4.543.333 & 7.311.667 \\
\bottomrule
\end{tabularx}
\end{table}

\textbf{``A partir de quando se tem lucro?''} --- três respostas, da mais simples à mais rigorosa (modelo C):
\begin{enumerate}
\item \textbf{Break-even operacional} (fluxo anual deixa de ser negativo): \textbf{Ano 2} (fluxo $+$\rs~88 mil).
\item \textbf{Payback simples} (recupera o capital, sem descontar): $\approx$ \textbf{2,7 anos} (o acumulado cruza zero no Ano 3).
\item \textbf{Payback descontado} (recupera já pagando a TMA --- definição mais rigorosa de início de criação de riqueza): $\approx$ \textbf{3,1 anos}. Como o horizonte é de 5 anos, o projeto é \textbf{viável}: VPL > 0, TIR $\gg$ TMA e IBC > 1.
\end{enumerate}

\textbf{Análise de sensibilidade (modelo C --- o motor é a adoção):}

\begin{table}[h]\small\centering
\begin{tabularx}{\textwidth}{l X r c X}
\toprule
\rowcolor{accent!12} Cenário & Premissa & VPL (20\%) & TIR & Veredito \\
\midrule
Otimista & adoção $+$25\%, preço $+$10\% & \rs~6.048.934 & 133,8\% & Muito viável \\
Base & premissas centrais & \rs~2.937.891 & 85,6\% & Viável \\
Pessimista & adoção $-$40\%, preço $-$15\% & $-$\rs~381.183 & 5,1\% & Inviável no horizonte \\
\bottomrule
\end{tabularx}
\end{table}

A viabilidade \textbf{depende da adoção}: se a contratação de hectares ficar $\approx$40\% abaixo do esperado, o projeto não se paga em 5 anos. A tese é, no fundo, sobre \textbf{velocidade de aquisição de hectares}.

\section{Case Real (dados a coletar)}
Para fechar o trabalho, substituir as premissas por dados verificáveis: (1) tamanho de mercado atual com fonte citada e datada; (2) preço e escopo de concorrentes reais (Perfect Flight, SIMA, SSCrop); (3) caso de uma fazenda representativa do Centro-Oeste --- área, cultura, \textbf{perda média por safra}, gasto com insumos, horas de inspeção --- traduzido em \textbf{\rs~economizados/ano} (a proposta de valor quantificada que justifica o preço/ha); (4) cotações reais de drone, sensores, nuvem e equipe para o CAPEX/OPEX; e (5) recálculo dos indicadores na planilha. Como a equipe está em Goiânia, no coração do agro, o acesso a um caso real é uma vantagem.

\section{Riscos e mitigações}
\begin{itemize}
\item \textbf{Adoção mais lenta que o previsto (risco nº 1):} pilotos que provem o ROI em \rs, canal via cooperativas e revendas, contratos plurianuais.
\item \textbf{Conectividade no campo:} mitigada por design (operação offline-first no gateway de borda).
\item \textbf{Custo de nuvem/processamento:} inferência leve na borda reduz o OPEX variável.
\item \textbf{Qualidade do dado e segurança:} validação humana (agrônomo) sobre as recomendações da IA antes da ação em campo.
\item \textbf{Concentração de receita:} diversificar a base para reduzir dependência de poucos clientes grandes.
\end{itemize}

\section{Conclusões e próximos passos}
O problema é real e validado pelo mercado; a solução tem diferencial defensável (integração + offline-first + prescrição acionável); é construível em fases. No modelo recomendado (Híbrido), a análise indica \textbf{viabilidade}: VPL positivo e TIR muito acima da TMA em todos os níveis testados (14,5\% / 20\% / 30\%), com \textbf{payback descontado de $\approx$3 anos}. A honestidade da análise está na sensibilidade: se a adoção decepcionar ($\approx$40\% abaixo), o projeto \textbf{não se paga no horizonte} --- logo, a decisão de investir é uma decisão sobre a capacidade de \textbf{acelerar a adoção de hectares}.

\textbf{Próximos passos (MVP até julho):} (1) construir o ciclo ``voo $\to$ análise de IA $\to$ alerta'' com uma fonte de dado e um caso de uso; (2) rodar um piloto em uma fazenda real e medir o \rs~economizado; (3) coletar cotações e preços de concorrentes para recalibrar a planilha; (4) consolidar os entregáveis (apresentação, vídeo, artigo) sobre estes números.

\section{Referências e fontes a verificar}
As premissas devem ser confirmadas, com fonte citada e datada, antes da entrega final:
\begin{itemize}
\item \textbf{TMA / Selic:} Banco Central do Brasil --- decisão do Copom (jun/2026). Reconfirmar a taxa vigente na data da apresentação.
\item \textbf{Concorrentes:} Perfect Flight, SIMA e SSCrop --- páginas oficiais das empresas (verificar os números de hectares e o escopo de cada serviço).
\item \textbf{Conectividade rural:} usar uma fonte oficial (ex.: pesquisa setorial de conectividade no campo) e citar o ano da medição.
\item \textbf{CAPEX/OPEX e preços por hectare:} são premissas próprias deste estudo, sem fonte externa; substituir por cotações reais de fornecedores e pela tabela de preços dos concorrentes.
\end{itemize}

\vspace{4pt}
\noindent\colorbox{accent!8}{\parbox{\dimexpr\textwidth-2\fboxsep}{\small\itshape \textbf{Limite de confiança:} os números financeiros são premissas conferidas, não dados verificados. A defesa final depende de substituí-los pelos dados reais da Seção 4.}}

\end{document}
